\hwhat{Exploratory round 1b (non-holdout, corrected activity table, leading-@ targets, no future-kick isolation, day fixed effect): a nudge adds 1.16 [0.75, 1.57] active minutes to its named target over 30 min in \#51 (first nudges 1.38, repeats 0.90; regime-III 4 h days 0.93), bystanders $\approx0$. The response starts at the target's next model call (median read-out 108 s; read-out-aligned onset 1 min): round 1's 4-min dead time was dropped events plus the isolation rule. Human messages no longer move activity. FDT fails in shape; kicked activations persist ($X\approx1.9$); mean-field Glauber underpredicts decay times. Nudger off (NE43): $n$ unchanged, traps not longer. \textbf{Locked holdout NE21} (hours $4\to8\to4\to8$ h): \textbf{C1 falsified in reverse} ($n$ 0.31, 0.67, 0.46, 0.89); C1 does not use the corrected table and stands.}

\hmodel{Models 02 (linear response, fluctuation--dissipation) and 09 (Hawkes). $n_i(m)\in\{0,1\}$: agent $i$ active in minute $m$. A kick (nudge, human message) at $m_0$ is a field pulse; $G$: matched response; $r$: read-out delay to the target's next model call; $C$: activity autocorrelation; $\chi$: integrated $G$; $X$: FD ratio. Hawkes: $\mu_d b(t)$ daily baseline, $\phi$ talk-triggered kernel, $\eta_c$ messages triggered per class-$c$ kick, $n$ branching ratio. Mean field: loop gain $K$, agent time $\tau_0$.}

\hmath{\vspace{-5pt}\begin{gather*}
G(\tau)=\mathbb E[n_i(m_0{+}\tau)\mid\text{kick}]-\mathbb E[n_i(m_0{+}\tau)\mid\text{control}]-\delta_d(\tau)\\
G(\tau)\approx\textstyle\sum_r p(r)\,\Delta R(\tau-r),\qquad \text{FDT: }\chi(\tau)=X\,[C(0)-C(\tau)]\\
\lambda(t)=\mu_d b(t)+\textstyle\sum_k\phi(t-t_k)+\sum_c\eta_c\sum_j\gamma_c e^{-\gamma_c(t-t_j)}\\
n=\textstyle\int_0^\infty\phi(\tau)\,d\tau,\qquad \tau=\tau_0/(1-K)
\end{gather*}\vspace{-7pt}}

\hobs{../figures/r1b_summary_obs.pdf}{(a) A nudged agent's activity against matched controls (corrected design, 95\% day-block band), aligned on the nudge and on the target's receiving call. Realigned, the response is immediate: the delay seen from the nudge is the wait until the target's next call. (b) Branching ratio per hours segment on the holdout; every switch went the other way, at the edge of week-to-week noise.}

\htelos{The useful part is plain and now sharper. A nudge works on the agent it names and nobody else, at that agent's next model call, so its lag is the target's pause schedule: a nudge to an agent paused for 10 minutes does nothing until it wakes. A first nudge buys about 1.4 active minutes, a repeat less. Human messages change what agents say, not how much they work. The physics vocabulary (Green's functions, FDT, effective temperature) mostly shows that a field-on-activity model is wrong. The one reversal test, longer days lowering self-excitation, failed in the opposite direction; switching the nudger off left self-excitation unchanged.}

\hbottom{A nudge acts only on its target, at the target's next model call (the apparent minutes-long dead time was a measurement artifact), and no hours or nudger switch lowered the swarm's self-excitation as predicted.}

\hresults{\scriptsize\setlength{\tabcolsep}{2pt}\begin{tabular}{@{}p{0.33\columnwidth}p{0.47\columnwidth}p{0.17\columnwidth}@{}}
\textbf{Prediction} & \textbf{Round 1 $\to$ round 1b} & \textbf{Verdict (1b)}\\\hline
P1 nudge: A60 $>0$, peak $\le$5 min & peak $\sim$15 min $\to$ \#51 peak 4, A60 1.65 [0.79, 2.52] & yes (\#51, 4 h)\\
Nudge A30 (isolated $\to$ all nudges) & 1.54 $\to$ 1.16 (\#51), 0.93 (4 h); isolation inflates $\times$1.3--1.9 & corrected\\
First vs repeat nudge (\#51) & -- $\to$ 1.38 vs 0.90 & sub-additive\\
P2 bystanders $\le$0.2 of target & yes $\to$ 0.03 (\#51) & yes\\
P3 human messages move activity & mentioned 5.6 $\to$ 0.6 (n.s.); room $\approx0$ & no\\
P6 FDT violated, nudge $X<1$ & violated, $X$ 1.54 $\to$ 1.88 & violated; dir.\ no\\
MF-P1/P3 ($K$ subcritical, self-consistent) & $K$ III 0.26 $\to$ 0.32 & yes\\
NE43 nudger off: $n$ same, traps longer & $\Delta n$ +0.06; inactive runs +0.5\% & mixed\\
NE44 read-out moves with pause default & 43 s vs 108 s; onset lag equal from read-out & mixed\\
C1 (holdout) $n$ lower on 8 h days & 0/3 switches; not affected by 1b & falsified\\
\end{tabular}}

\hobsb{../figures/r1b_summary_obs2.pdf}{(a) Cumulative response of \#51 nudge targets by how long the nudge waited to be read: read within 3 min, the response starts at once; read after 10 min (long pauses), the target stays below its controls. (b) NE43: Hawkes $n$ with the daily bookends and nudges, with nudges only, and with neither.}

\hcaveats{The \#51 pre-window placebo is +0.25 [0.05, 0.46]: some selection by the nudger's text trigger remains, so the level may be up to a quarter too high. Long pauses are not in the matching strata. Regime-I nudges (NE10) fail the placebo. Hawkes $n$ may absorb within-day nonstationarity. The holdout's C2, C4 and MF-C used the buggy activity table and the isolation rule; re-running them is Vivian's call.}

\hnext{Add pause state to the matching strata; fit a read-out + transient kernel.}
